\chapter{ВАЛИДАЦИЯ, РЕЗУЛЬТАТЫ ЭКСПЕРИМЕНТОВ И АНАЛИЗ}
\label{ch:results}
\sloppy


\section{Цель и схема экспериментальной валидации}
\label{sec:4-1}

Постановка задачи и метод сценарной оценки риска перегрузки залов сформулированы в главе~\ref{ch:method}; реализующий их программный комплекс описан в главе~\ref{ch:implementation}. Настоящая глава предъявляет результаты применения комплекса к программе конференции Mobius 2025 Autumn и содержит экспериментальную валидацию принятых архитектурных решений.

Валидация в работе понимается в смысле, принятом в литературе по имитационному моделированию: она не утверждает, что построенная модель воспроизводит фактическое поведение аудитории конкретной конференции, поскольку для такого утверждения необходимы данные о фактической посещаемости, в постановке работы недоступные. Валидация состоит в проверке трёх свойств комплекса. Первое: модель сохраняет обязательные граничные свойства --- при свободной вместимости риск перегрузки исчезает, при сжатой растёт; при нулевом весе компонента следования рекомендации политики неразличимы по показателям загрузки, при единичном весе различимы максимально. Второе: относительные сравнения политик и вариантов программы остаются устойчивыми при варьировании параметрических осей. Третье: проверяется степень согласованности ранжирования политик в параметрическом и LLM-агентском симуляторах на общем подмножестве конфигураций.

Все три свойства проверяются на едином плане эксперимента: пятидесяти точках латинского гиперкуба с тремя независимыми случайными зёрнами на каждую комбинацию <<точка $\times$ политика>>, репрезентативном подмножестве из двенадцати точек для агентского симулятора и фиксированных значениях параметров стохастичности и базовой вероятности отказа от посещения. Полный объём прогонов составляет 486 единичных итераций параметрического симулятора и 48 единичных итераций LLM-агентского симулятора.

Структура главы соответствует логической последовательности проверок. Подраздел~\ref{sec:4-2} предъявляет результаты граничной верификации модели как блокирующего фильтра перед содержательными выводами. Подраздел~\ref{sec:4-3} описывает результаты полного параметрического прогона --- общее распределение показателей по политикам и подмножество конфигураций, в которых риск перегрузки физически возникает. Подраздел~\ref{sec:4-4} раскрывает попарные сравнения политик и устойчивость ранжирования. Подраздел~\ref{sec:4-5} рассматривает эффект локальных модификаций программы. Подраздел~\ref{sec:4-6} исследует чувствительность результата к параметрическим осям. Подраздел~\ref{sec:4-7} посвящён согласованности двух симуляторов. Подраздел~\ref{sec:4-8} фиксирует ограничения эксперимента и угрозы валидности. Подраздел~\ref{sec:4-9} подводит итог.

Все численные значения, приведённые в главе, получены из зафиксированных артефактов полного экспериментального прогона на программе Mobius 2025 Autumn. Сводные результаты представлены в таблицах~\ref{tab:ec-verification}--\ref{tab:cross-validation}; состав артефактов эксперимента приведён в приложении~\ref{app:artefacts}. Конкретные технические детали реализации не повторяются: они описаны в главе~\ref{ch:implementation}.


\section{Верификация граничных свойств модели}
\label{sec:4-2}

Граничная верификация модели --- блокирующий фильтр: до её прохождения сравнения политик и вариантов программы не интерпретируются содержательно. Четыре граничных свойства из подраздела~2.4.4 главы~\ref{ch:method} проверяются автоматизированными тестами на минимально достаточных конфигурациях с усреднением по нескольким случайным зёрнам.

Свойство первое. При сильно избыточной номинальной вместимости --- множитель не менее 3.0 относительно номинальной --- среднее относительное превышение вместимости тождественно равно нулю независимо от выбранной политики. Это базовая внутренняя согласованность: при заведомо достаточной вместимости перегрузка не может возникнуть ни при какой стратегии распределения участников.

Свойство второе. При уменьшении номинальной вместимости показатели риска перегрузки не убывают. Проверяется на сетке множителей вместимости с усреднением по пяти независимым случайным зёрнам. Тест считается пройденным, если среднее значение показателя на каждом следующем (более сжатом) уровне не ниже значения на предыдущем (более свободном) уровне.

Свойство третье. При нулевом весе компонента следования рекомендации различие политик по показателям загрузки исчезает. Реализованная функция полезности при нулевом весе не зависит от рекомендации, и распределение выбора участника определяется только содержательной близостью и компонентом социального заражения. Этот инвариант контролируется тестом на пословное совпадение протоколов прогона при разных политиках при одинаковых случайных зёрнах: при нулевом весе все политики порождают один и тот же протокол с точностью до записи сигналов о рекомендациях.

Свойство четвёртое. При единичном весе компонента следования рекомендации различие политик по показателям загрузки максимально. Проверяется на стрессовой конфигурации вместимости: размах значений показателя между политиками в этом режиме существенно превосходит коэффициент вариации того же показателя в режиме нулевого веса.

В составе комплекса реализованы десять автоматизированных тестов: четыре теста соответствуют перечисленным свойствам, остальные шесть расширяют их на конфигурации с ненулевым весом компонента социального заражения. Сводный результат граничной верификации представлен в таблице~\ref{tab:ec-verification}.

\begin{small}
\begin{longtable}{|>{\raggedright\arraybackslash}p{3.2cm}|>{\raggedright\arraybackslash}p{4.2cm}|>{\raggedright\arraybackslash}p{5.5cm}|>{\centering\arraybackslash}p{1.6cm}|}
\caption{Результаты граничной верификации модели}\label{tab:ec-verification}\\
\hline
\textbf{Проверка} & \textbf{Проверяемое свойство} & \textbf{Ожидаемое поведение} & \textbf{Результат} \\
\hline
\endfirsthead
\caption*{Продолжение таблицы~\ref{tab:ec-verification}}\\
\hline
\textbf{Проверка} & \textbf{Проверяемое свойство} & \textbf{Ожидаемое поведение} & \textbf{Результат} \\
\hline
\endhead
\hline
\endfoot
\hline
\endlastfoot
EC1 & Свободная вместимость & При множителе вместимости $\geq 3.0$ среднее относительное превышение вместимости равно нулю для всех политик & Пройдено \\
\hline
EC2 & Монотонность по вместимости & При уменьшении множителя вместимости средний риск не убывает (усреднение по 5 зёрнам) & Пройдено \\
\hline
EC3 & Граница по $w_{rec}$ снизу & При $w_{rec} = 0$ протоколы прогонов разных политик пословно совпадают & Пройдено \\
\hline
EC4 & Граница по $w_{rec}$ сверху & При $w_{rec} = 1$ размах показателя между политиками существенно превосходит шум репликации & Пройдено \\
\hline
EC1 (расширение) & EC1 при $w_{gossip} = 0.3$ & Превышение вместимости остаётся нулевым при включённом социальном заражении и свободной вместимости & Пройдено \\
\hline
EC2 (расширение) & EC2 при $w_{gossip} = 0.3$ & Монотонность по вместимости сохраняется при включённом социальном заражении & Пройдено \\
\hline
EC3 (расширение) & EC3 при $w_{gossip} > 0$ & Совпадение протоколов прогонов сохраняется при включённом социальном заражении & Пройдено \\
\hline
EC4 (расширение) & EC4 при ненулевом $w_{gossip}$ & На границе симплекса при $w_{rec} = 0.7$ и $w_{gossip} = 0.3$ различие политик по показателям загрузки сохраняется & Пройдено \\
\hline
Сохранение эффекта политики с учётом загрузки & Сохранение эффекта политики с учётом загрузки залов при $w_{gossip} > 0$ & При $w_{gossip} = 0.3$ и стрессовой вместимости политика с учётом загрузки залов не уступает политике по релевантности по риску перегрузки & Пройдено \\
\hline
Монотонность концентрации & Усиление концентрации выбора по $w_{gossip}$ & Коэффициент Джини загрузки залов не убывает по сетке $w_{gossip} \in \{0,\ 0.3,\ 0.7\}$ & Пройдено \\
\end{longtable}
\end{small}

Граничная верификация пройдена полностью: 10 из 10 автоматизированных тестов завершились без отказов. Поэтому дальнейшие сравнения политик выполняются на реализации, удовлетворяющей обязательным граничным свойствам модели.

Прохождение граничной верификации означает, что реализованная функция полезности и схема разделения источников случайности удерживают разделение каналов влияния, заявленное в постановке: без этого условия попарные сравнения политик не имели бы методического смысла, поскольку политики не были бы инструментально отделены от состояния системы. Кроме того, модель не порождает контринтуитивных артефактов на границах параметрических осей. Содержательные сравнения политик и вариантов программы в остальных подразделах главы опираются именно на эту проверенную реализацию.

Граничная верификация не утверждает, что модель воспроизводит реальное поведение аудитории; такое утверждение при отсутствии данных о фактической посещаемости невозможно по построению. Она утверждает только, что обязательные свойства, без которых сценарный анализ был бы некорректен, в реализации не нарушены.


\section{Результаты параметрического сценарного эксперимента}
\label{sec:4-3}

Полный параметрический прогон выполнен на программе Mobius 2025 Autumn по плану, описанному в подразделе~\ref{sec:3-6}. Внутри каждой из пятидесяти точек гиперкуба перебирались политики с тремя независимыми случайными зёрнами; политика на основе языковой модели в параметрическом режиме запускалась только на двенадцати точках репрезентативного подмножества по соображениям стоимости. Суммарный объём --- 486 единичных итераций. Главный аналитический интерес представляет подмножество тех точек гиперкуба, в которых риск перегрузки фактически возникает: на остальных точках любое сравнение политик вырождается, поскольку при свободной вместимости все политики дают нулевой риск независимо от логики ранжирования.

Распределение центральных показателей по политикам в полном прогоне на пятидесяти точках представлено в таблице~\ref{tab:central-metrics}.

\begin{table}[h!]
\caption{Центральные показатели параметрического прогона по политикам ($n = 50$ точек гиперкуба, 3 реплики на точку)}
\label{tab:central-metrics}
\centering
\small
\begin{tabular}{|l|r|r|r|r|}
\hline
\textbf{Политика} & \textbf{Среднее превышение} & \textbf{Медиана} & \textbf{P75} & \textbf{Средняя релевантность} \\
\hline
Контрольная               & 0.0395 & 0.0000 & 0.0000 & 0.7308 \\
\hline
По релевантности          & 0.0449 & 0.0000 & 0.0066 & 0.7321 \\
\hline
С учётом загрузки залов   & 0.0336 & 0.0000 & 0.0000 & 0.7314 \\
\hline
\end{tabular}
\end{table}

Сравнение средних в таблице~\ref{tab:central-metrics} даёт первое впечатление о близости трёх политик: различия между ними лежат в пределах сотых долей единицы относительного превышения, а медианы у всех трёх политик равны нулю. Устойчивый сравнительный разрыв между политиками выявляется не средними по выборке, а попарным анализом внутри точки гиперкуба (подраздел~\ref{sec:4-4}), и сосредоточен на risk-positive подмножестве, описанном ниже. Тем не менее уже на уровне сводных характеристик распределения видно: политика с учётом загрузки залов даёт наименьшее среднее относительное превышение вместимости при практически неизменной средней релевантности выбранных докладов. Различие в средней релевантности между тремя политиками не превышает 0.002, то есть переход от чистого ранжирования по релевантности к учёту состояния залов не сопровождается заметной потерей содержательного качества рекомендаций.

Нулевая медиана риска перегрузки в полном прогоне --- это структурное свойство сценарного анализа, а не отсутствие эффекта политик: латинский гиперкуб по построению покрывает всё пространство параметрических конфигураций, в том числе режимы свободной вместимости, в которых риск нулевой по физическим причинам. Из пятидесяти точек гиперкуба ненулевой риск хотя бы у одной из трёх политик контрольного состава возникает в тринадцати. Сводка по этому подмножеству приведена в таблице~\ref{tab:risk-positive}.

\begin{table}[h!]
\caption{Подмножество точек с ненулевым риском перегрузки (risk-positive)}
\label{tab:risk-positive}
\centering
\small
\begin{tabular}{|>{\raggedright\arraybackslash}p{10cm}|r|}
\hline
\textbf{Показатель} & \textbf{Значение} \\
\hline
Всего точек LHS в плане эксперимента & 50 \\
\hline
Точек, на которых хотя бы у одной из политик П1--П3 возникает ненулевой риск & 13 \\
\hline
Доля risk-positive точек & 26~\% \\
\hline
Точек с ненулевым риском у политики по релевантности & 13 / 13 \\
\hline
Точек с ненулевым риском у контрольной политики & 11 / 13 \\
\hline
Точек с ненулевым риском у политики с учётом загрузки залов & 10 / 13 \\
\hline
Точек с ненулевым риском у политики на основе языковой модели (на 12 maximin) & 2 / 12 \\
\hline
\end{tabular}
\end{table}

Содержательный анализ риска должен в первую очередь смотреть на risk-positive подмножество: на остальных точках все политики дают нулевой риск по физическим причинам, и сравнение политик там вырождается. Различие частот ненулевого риска между контрольной политикой и политикой по релевантности отражает структурный эффект ранжирования по содержательной близости: при отсутствии учёта состояния залов рекомендация по релевантности концентрирует поток участников на одних и тех же темах, что в части конфигураций приводит к перегрузке одного из залов.

Среди risk-positive точек выделяется подмножество критически непреодолимых конфигураций --- точек, в которых размер аудитории превышает суммарную вместимость залов в параллельном слоте при текущем множителе вместимости. Таких конфигураций три: при множителях вместимости 0.534, 0.963 и 0.774 в сочетании с размером аудитории сто. Физическое переполнение в них неустранимо никакой стратегией распределения участников. Точки оставлены в анализе для прозрачности: они задают границу применимости сценарного сравнения и позволяют организатору отделить случаи, в которых рекомендация бессильна, от случаев, в которых политика даёт выигрыш.

Полная таблица единичных итераций сохранена в формате long-table с явным указанием значений всех осей конфигурации, идентификатора политики, номера реплики и значений всех показателей групп, описанных в подразделе~2.2.3. Эта таблица служит исходным материалом для всех последующих процедур постобработки, описанных в подразделах~\ref{sec:4-4}--\ref{sec:4-7}.


\section{Попарное сравнение политик и устойчивость выводов}
\label{sec:4-4}

Сравнение политик в работе ведётся попарно внутри одной точки гиперкуба, а не средним показателем по всему пространству конфигураций. Этот выбор содержательно обусловлен тем, что среднее по пространству конфигураций скрывает сценарную чувствительность результата, а в задаче с глубокой неопределённостью именно чувствительность представляет интерес.

Для каждой пары политик в каждой точке гиперкуба сравниваются три значения центральных показателей --- по три независимых случайных зёрна на политику --- и фиксируется, какая политика лидирует. Сравнение выполняется в двух режимах: строгое неравенство и неравенство с допуском. Допуск выбран равным 0.005 для показателей риска и неравенства загрузки и 0.001 для показателя релевантности; эти значения выбраны на этапе постобработки и зафиксированы в таблице $\varepsilon$. Введение допуска позволяет различить случаи фактического превосходства одной политики над другой и случаи практически равных результатов в пределах статистического шума репликации. Доли точек, в которых политика выигрывает строго, выигрывает с допуском, не отличается с допуском, проигрывает строго и проигрывает с допуском, в сумме дают единицу с точностью до округления.

Сводный результат попарного сравнения трёх контрольных политик в полном прогоне для метрики среднего относительного превышения вместимости при допуске $\varepsilon = 0.005$ приведён в таблице~\ref{tab:pairwise-overload}.

\begin{table}[h!]
\caption{Попарное сравнение политик по риску перегрузки в полном прогоне (50 точек, $\varepsilon = 0.005$)}
\label{tab:pairwise-overload}
\centering
\small
\begin{tabularx}{\linewidth}{|>{\raggedright\arraybackslash}p{3.6cm}|>{\centering\arraybackslash}p{1.7cm}|>{\raggedright\arraybackslash}p{3.0cm}|>{\centering\arraybackslash}p{1.5cm}|>{\centering\arraybackslash}p{1.5cm}|>{\raggedright\arraybackslash}X|}
\hline
\textbf{Пара политик} & \textbf{Равны с $\varepsilon$} & \textbf{Победитель за пределами допуска} & \textbf{Доля побед} & \textbf{Доля побед второй} & \textbf{Вывод} \\
\hline
Контрольная vs по релевантности & 86~\% & Контрольная & 14~\% & 0~\% & Контрольная превосходит политику по релевантности в 14~\% точек; обратная победа за пределами допуска не зафиксирована \\
\hline
Контрольная vs с учётом загрузки залов & 86~\% & С учётом загрузки залов & 14~\% & 0~\% & Политика с учётом загрузки залов превосходит контрольную политику в 14~\% точек; обратная победа за пределами допуска не зафиксирована \\
\hline
По релевантности vs с учётом загрузки залов & 78~\% & С учётом загрузки залов & 22~\% & 0~\% & Политика с учётом загрузки залов превосходит политику по релевантности в 22~\% точек; обратная победа за пределами допуска не зафиксирована \\
\hline
\end{tabularx}
\end{table}

Наиболее устойчивый результат попарного сравнения: политика по релевантности не выигрывает у политики с учётом загрузки залов ни в одной из 50 точек гиперкуба ни строго (\texttt{win\_strict = 0}), ни за пределами допуска $\varepsilon = 0.005$ (\texttt{win\_eps = 0}). Политика с учётом загрузки залов, напротив, выигрывает у политики по релевантности в 26~\% точек строго и в 22~\% точек за пределами допуска. На risk-positive подмножестве из тринадцати точек политика с учётом загрузки залов не уступает максимуму из контрольной политики и политики по релевантности на всех 13 / 13 точках в пределах допуска и строго снижает риск перегрузки относительно этого максимума в 11 / 13 точек ($\approx$~85~\%).

Шесть наибольших снижений риска при переходе от политики по релевантности к политике с учётом загрузки залов представлены в таблице~\ref{tab:risk-reductions}.

\begin{table}[h!]
\caption{Наибольшие снижения риска при переходе от политики по релевантности к политике с учётом загрузки залов (risk-positive подмножество)}
\label{tab:risk-reductions}
\centering
\small
\begin{tabular}{|r|r|r|r|r|r|}
\hline
\textbf{Точка LHS} & \textbf{Множитель} & \textbf{Аудитория} & \textbf{Риск (П2)} & \textbf{Риск (П3)} & \textbf{Снижение} \\
\hline
49 & 0.774 & 100 & 0.5449 & 0.3558 & 0.1891 \\
\hline
26 & 0.629 &  60 & 0.1706 & 0.0040 & 0.1667 \\
\hline
35 & 0.963 & 100 & 0.0606 & 0.0025 & 0.0581 \\
\hline
 3 & 0.534 & 100 & 1.2222 & 1.1759 & 0.0463 \\
\hline
36 & 0.696 &  60 & 0.1424 & 0.1111 & 0.0312 \\
\hline
18 & 1.040 &  60 & 0.0214 & 0.0000 & 0.0214 \\
\hline
\end{tabular}
\end{table}

Самое крупное снижение зарегистрировано на точке 49: одна из трёх критически непреодолимых конфигураций (audience > effective per-slot capacity), на которой политика с учётом загрузки залов снижает риск с 0.5449 до 0.3558 при том, что физическое переполнение полностью неустранимо. Точка 26 показывает крупное абсолютное снижение риска: с 0.1706 до 0.0040, то есть почти до нулевого значения.

Устойчивость результата отдельной политики на отдельной точке проверяется средствами анализа разброса показателей между репликами одной комбинации <<точка $\times$ политика>>. На полном прогоне трёх контрольных политик зафиксированы 24 сочетания <<точка $\times$ политика $\times$ показатель>> из 600, в которых отношение стандартного отклонения по репликам к абсолютному значению среднего превышает 0.5; на двенадцати точках с политикой на основе языковой модели --- 4 из 36. Все эти сочетания соответствуют режимам редких событий перегрузки, в которых редкая ненулевая реплика порождает большое относительное стандартное отклонение при малом абсолютном значении среднего; центральные выводы попарного сравнения политик в подразделах~\ref{sec:4-3}--\ref{sec:4-4} не зависят от этих конкретных комбинаций. Подробное распределение разброса по репликам сохранено в составе артефактов постобработки (приложение~\ref{app:artefacts}).

В подмножестве из двенадцати точек, на которых запускалась также политика на основе языковой модели, попарные сравнения по показателю риска перегрузки согласуются с полным прогоном. Политика с учётом загрузки залов даёт одинаковые с допуском результаты с контрольной политикой в 83~\% точек и выигрывает за пределами допуска в 17~\%; с политикой по релевантности --- на тех же 83~\% одинаковые с допуском и в 17~\% точек выигрывает за пределами допуска. По показателю среднего относительного превышения вместимости политика с учётом загрузки залов даёт минимальное среднее (0.0093), политика на основе языковой модели --- 0.0109, контрольная политика --- 0.0127, политика по релевантности --- 0.0136. Эти средние не сопоставимы со средними по пятидесяти точкам: десять из двенадцати точек подмножества попадают в режимы свободной вместимости, и абсолютные значения смещены вниз. Политика на основе языковой модели присутствует только в этом подмножестве.

Тепловая карта ранжирования политик по четырём показателям на двенадцати точках репрезентативного подмножества приведена на рисунке~\ref{fig:ranking-heatmap}. Цветовая шкала отражает позицию политики в ранжировании: тёмные клетки соответствуют верхним позициям, светлые --- нижним. Видна устойчивость лидерства политики с учётом загрузки залов на большинстве точек по показателям риска при разреженной структуре ранжирований по показателям, на которых параметрический симулятор вырождается.

\begin{figure}[h!]
\centering
\includegraphics[width=0.85\linewidth]{analysis_ranking_heatmap_maximin.png}
\caption{Ранжирование политик по четырём показателям на 12 точках maximin-подмножества}
\label{fig:ranking-heatmap}
\end{figure}


\section{Эффект локальных модификаций программы}
\label{sec:4-5}

Оператор локальных модификаций программы, описанный в подразделе~\ref{sec:3-4}, генерирует допустимые варианты программы перестановками пар докладов между разными временными слотами одного дня. Индекс варианта программы фигурирует как одна из шести осей параметрической конфигурации; в полном прогоне на пятидесяти точках индекс варианта принимает шесть уровней, соответствующих исходной программе и пяти её локальным модификациям. Поскольку выборка по плану эксперимента строится латинским гиперкубом, разные уровни индекса варианта программы попадают в разные точки гиперкуба, не выровненные между собой по остальным пяти осям. Это техническое свойство выборки имеет принципиальное методическое следствие: попарные сравнения <<исходная программа против модификации $k$>> при фиксированных значениях остальных осей в текущей выборке не выполнены, и любые агрегаты по уровню индекса варианта программы на полной выборке носят диагностический характер, а не утверждают причинной разности между уровнями.

Диагностические агрегаты среднего относительного превышения вместимости по уровням индекса варианта программы в полном прогоне для политики по релевантности приведены в таблице~\ref{tab:program-effect}.

\begin{table}[h!]
\caption{Диагностические агрегаты по индексу варианта программы (политика по релевантности)}
\label{tab:program-effect}
\centering
\small
\begin{tabularx}{\linewidth}{|r|r|r|>{\raggedright\arraybackslash}X|}
\hline
\textbf{Индекс варианта} & \textbf{Число точек LHS} & \textbf{Среднее превышение} & \textbf{Интерпретация} \\
\hline
0 (исходная программа) &  7 & 0.0099 & Контрольный уровень \\
\hline
1 &  9 & 0.0211 & Заметное изменение в большую сторону \\
\hline
2 &  9 & 0.0633 & Заметное изменение в большую сторону \\
\hline
3 &  8 & 0.0000 & Все точки данного уровня в режиме свободной вместимости \\
\hline
4 &  7 & 0.1762 & Наибольшее агрегатное превышение в выборке \\
\hline
5 & 10 & 0.0183 & Близко к контрольному уровню \\
\hline
\end{tabularx}
\end{table}

Таблица не доказывает причинный эффект конкретной перестановки: уровни индекса программы в плане эксперимента не сопоставлены по остальным осям LHS. Знаковый тест разности между уровнями 0 и 4 на cross-product парах даёт p-значение около 0.84; гипотеза об отсутствии разности на текущей выборке не отвергается. Размах агрегата между уровнями 3 и 4 (от 0.0000 до 0.1762) отражает прежде всего различие остальных осей вокруг разных уровней индекса программы, а не различие самих модификаций.

Результат раздела состоит в том, что оператор $\Phi$ интегрирован в сценарный эксперимент как полноценная параметрическая ось наравне с вместимостью, размером аудитории и весами компонентов модели поведения; сводки по уровням этой оси (таблица~\ref{tab:program-effect}) дают сценарные характеристики распределения показателей риска по вариантам программы. Положение №4 защищаемых результатов закрывается именно как количественные сценарные оценки распределения по оси конфигурации, а не как причинная разность между конкретными перестановками. Сравнение <<исходная программа против модификации~$k$>> при фиксированных остальных осях относится к локальной оптимизации программы и оставлено за рамками работы; этот шаг зафиксирован в подразделе~\ref{sec:4-8} как направление развития.

Вторая сторона того же результата --- реализованный комплекс позволяет численно сопоставить эффект изменения структуры программы с эффектом изменения политики рекомендаций в одной экспериментальной обстановке. Эта возможность отвечает четвёртому защищаемому положению работы о совместном оценивании вклада алгоритмической логики и структуры программы в показатели риска.


\section{Анализ влияния параметрических осей}
\label{sec:4-6}

Чувствительность результатов к параметрическим осям конфигурации исследуется через сопоставление подмножеств точек гиперкуба, отличающихся значениями одной из осей при варьируемых остальных. Такое сопоставление конфаундировано: варьирующиеся остальные оси вносят собственный вклад в наблюдаемую разность, поэтому каждая сводка трактуется как диагностический сигнал о направлении и порядке эффекта, а не как причинная оценка. Глобальные методы оценки чувствительности с разложением вклада осей рассматривались на этапе проектирования и отвергнуты как избыточные для задачи попарного сравнения политик внутри точки.

Влияние множителя вместимости и веса компонента следования рекомендации в данных согласуется со свойствами граничной верификации (подраздел~\ref{sec:4-2}): множитель меньше единицы концентрирует случаи ненулевого риска (10 из 50 точек), множитель от двух до трёх даёт нулевой риск во всех 19 точках; различие политик растёт по мере увеличения веса компонента следования рекомендации. Эти зависимости не приносят нового содержательного знания и приведены в составе сводок чувствительности (приложение~\ref{app:artefacts}).

Распределение среднего относительного превышения вместимости по бакетам веса социального заражения представлено в таблице~\ref{tab:gossip-buckets}.

\begin{table}[h!]
\caption{Средний риск перегрузки по бакетам веса социального заражения $w_{gossip}$}
\label{tab:gossip-buckets}
\centering
\small
\begin{tabular}{|c|r|r|r|r|}
\hline
\textbf{Бакет $w_{gossip}$} & \textbf{Число точек} & \textbf{Контрольная} & \textbf{По релевантности} & \textbf{С учётом загрузки залов} \\
\hline
$[0,\ 0.25)$  & 23 & 0.0061 & 0.0123 & 0.0003 \\
\hline
$[0.25,\ 0.5)$ & 19 & 0.0891 & 0.0957 & 0.0822 \\
\hline
$[0.5,\ 0.7]$ &  8 & 0.0178 & 0.0178 & 0.0139 \\
\hline
\end{tabular}
\end{table}

Графически тот же агрегат показан на рисунке~\ref{fig:gossip-bucket-bar}: высота столбцов отражает средний риск перегрузки по политикам в каждом из трёх бакетов веса социального заражения.

\begin{figure}[h!]
\centering
\includegraphics[width=0.85\linewidth]{analysis_gossip_bucket_bar.png}
\caption{Средний риск перегрузки по бакетам веса социального заражения}
\label{fig:gossip-bucket-bar}
\end{figure}

В данных наблюдается немонотонная связь между весом социального заражения и средним риском: максимум среднего риска приходится на средний бакет $w_{gossip}$. Это диагностический результат, а не причинная оценка: остальные оси LHS в бакетах не фиксированы, и наблюдаемый пик в среднем бакете может объясняться конфаундированностью с осями вместимости и размера аудитории, а не свойствами самого канала социального заражения. Внутри каждого бакета политика с учётом загрузки залов сохраняет минимальное среднее значение риска среди трёх политик контрольного состава.

Влияние размера аудитории и источника оценки популярности докладов в данных умеренное и не порождает новых содержательных закономерностей; соответствующие сводки сохранены в составе сводок чувствительности (приложение~\ref{app:artefacts}).

Связь риска и релевантности на пятидесяти точках исследуется через двумерное распределение пар <<риск перегрузки $\times$ релевантность>>. Сводка по частоте trade-off приведена в таблице~\ref{tab:tradeoff}.

\begin{table}[h!]
\caption{Частота trade-off между снижением риска и потерей релевантности ($\varepsilon_{\text{risk}} = 0.005$, $\varepsilon_{\text{utility}} = 0.001$)}
\label{tab:tradeoff}
\centering
\small
\begin{tabularx}{\linewidth}{|>{\raggedright\arraybackslash}X|r|r|r|}
\hline
\textbf{Подмножество} & \textbf{Всего комбинаций} & \textbf{С trade-off} & \textbf{Доля} \\
\hline
Полный прогон трёх основных политик & 150 & 11 & 7.3~\% \\
\hline
Подмножество с участием LLM-политики &  48 &  3 & 6.3~\% \\
\hline
\end{tabularx}
\end{table}

Двумерное распределение пар <<риск $\times$ релевантность>> по всем точкам и политикам показано на рисунке~\ref{fig:risk-utility-scatter}. Большинство точек сгруппированы в зоне нулевого риска при сопоставимой релевантности; точки trade-off лежат в правой верхней части картины и не образуют доминирующей моды.

\begin{figure}[h!]
\centering
\includegraphics[width=0.85\linewidth]{analysis_risk_utility_scatter.png}
\caption{Распределение пар <<риск перегрузки $\times$ релевантность>> по политикам}
\label{fig:risk-utility-scatter}
\end{figure}

Trade-off редок в обоих подмножествах: на большинстве точек гиперкуба политика, лидирующая по снижению риска, не сопровождает это лидерство заметной потерей содержательного качества рекомендаций. Это согласуется с близкими значениями средней релевантности у трёх политик в таблице~\ref{tab:central-metrics}.


\section{Согласованность параметрического и LLM-агентского симуляторов}
\label{sec:4-7}

Вторая проверка относительного порядка политик выполнена прогоном LLM-агентского симулятора на репрезентативном подмножестве из двенадцати точек гиперкуба, отобранных алгоритмом максимизации минимального попарного расстояния с принудительным включением точки исходной программы. Внутри одной точки все четыре политики прогоняются на одной и той же синтетической аудитории. Объём прогона --- 48 единичных итераций (двенадцать точек $\times$ четыре политики $\times$ одно случайное зерно). Технические характеристики прогона --- фактическая стоимость, время выполнения, число вызовов внешней языковой модели и доля ошибок парсинга --- сохранены в метаданных эксперимента и в основном тексте не используются.

Согласованность ранжирования политик между параметрическим и LLM-агентским симуляторами оценивается по четырём ключевым показателям. Для каждой из двенадцати точек подмножества вычисляются ранжирования политик в обоих симуляторах, и попарное соответствие ранжирований количественно характеризуется коэффициентом ранговой корреляции Спирмена. Минимальная медианная корреляция, при которой согласованность считается принятой, фиксирована на уровне 0.5. Сводный результат представлен в таблице~\ref{tab:cross-validation}.

\begin{table}[h!]
\caption{Согласованность ранжирования политик между параметрическим и LLM-агентским симуляторами на 12 точках подмножества}
\label{tab:cross-validation}
\centering
\small
\begin{tabularx}{\linewidth}{|>{\raggedright\arraybackslash}p{4.0cm}|>{\centering\arraybackslash}p{2.0cm}|>{\centering\arraybackslash}p{1.6cm}|>{\centering\arraybackslash}p{1.6cm}|>{\raggedright\arraybackslash}X|}
\hline
\textbf{Показатель} & \textbf{Невырожденных точек} & \textbf{Медиана $\rho$} & \textbf{Совпадение лидера} & \textbf{Интерпретация} \\
\hline
Средняя релевантность & 12 & 0.80 & 11 / 12 & Высокая согласованность \\
\hline
Доля слотов с переполнением & 2 & 0.74 & 2 / 2 & Диагностическая согласованность из-за малого числа невырожденных точек \\
\hline
Относительное превышение вместимости & 2 & 0.30 & 2 / 2 & Ниже порога $\rho \geq 0.5$, узкая выборка \\
\hline
Дисперсия загрузки залов в слоте & 12 & 0.40 & 11 / 12 & Ниже порога по $\rho$, но высокое совпадение лидера \\
\hline
\end{tabularx}
\end{table}

Графическая сводка коэффициентов ранговой корреляции Спирмена по метрикам перекрёстной проверки приведена на рисунке~\ref{fig:cross-validation-rho}.

\begin{figure}[h!]
\centering
\includegraphics[width=0.85\linewidth]{cross_validation_rho_per_metric.png}
\caption{Коэффициенты ранговой корреляции Спирмена по метрикам перекрёстной проверки}
\label{fig:cross-validation-rho}
\end{figure}

Объединённая медиана коэффициента корреляции по четырём показателям составляет 0.554 при пороге 0.5. Эта агрегатная величина сама по себе диагностически малоинформативна: она усредняет 28 значений с разной долей валидных точек --- по релевантности и по дисперсии загрузки по 12 точек, по показателям переполнения по 2 точки из-за вырождения параметрического ранжирования на 10 из 12 безопасных точек подмножества. Содержательную картину даёт построчная разбивка таблицы~\ref{tab:cross-validation}. Подробные данные по точкам с пометкой <<degenerate>> для каждой точки с вырожденным ранжированием сохранены в составе артефактов перекрёстной проверки (приложение~\ref{app:artefacts}).

Согласованность параметрического и LLM-агентского симуляторов наблюдается по части показателей: уверенно --- по средней релевантности (12 точек) и по верхнему лидерству политик в дисперсии загрузки (11 / 12 точек); по показателям переполнения корреляция оценивается на 2 невырожденных точках из 12 и остаётся диагностической. Интерпретация расхождений по метрикам переполнения ограничена структурой выборки: десять из двенадцати точек параметрически безопасны, поэтому ранжирование по этим метрикам вырождается. На текущем подмножестве нельзя отделить возможное расхождение моделей от эффекта малой доли невырожденных точек. Третьему защищаемому положению работы это отвечает в части реализации двух механизмов отклика и перекрёстной проверки на репрезентативном подмножестве.

Использованные в LLM-прогоне модели зафиксированы в метаданных эксперимента: политика на основе языковой модели работала на \texttt{openai/gpt-4o-mini}, агентский симулятор --- на \texttt{openai/gpt-5.4-nano}. Обе модели выбраны из соображений практической достижимости бюджета и не претендуют на максимально достижимое качество поведенческой модели; это сознательное ограничение, фиксируемое в подразделе~\ref{sec:4-8}.


\section{Ограничения эксперимента и угрозы валидности}
\label{sec:4-8}

Полученные результаты подлежат интерпретации в рамках принятой постановки задачи. Ограничения работы фиксируются как техническое условие задачи, не как защитная оговорка.

\textit{Отсутствие данных о фактической посещаемости.} Аудитория задана синтетическим порождающим процессом. Все численные значения показателей риска, загрузки и релевантности интерпретируются как сравнительные характеристики политик и вариантов программы внутри принятой модели, не как прогноз фактической посещаемости конкретной конференции.

\textit{Одна основная конференция.} Численные результаты полного прогона получены на программе Mobius 2025 Autumn. Перенос относительного порядка политик на программы других мероприятий требует отдельной экспериментальной проверки; универсальность выводов работы для произвольной программы конференции не утверждается.

\textit{Фиксированные параметры стохастичности и базовой вероятности отказа.} Параметр стохастичности softmax-распределения и базовая вероятность отказа от посещения слота в плане эксперимента не варьировались. Чувствительность результатов к этим параметрам в работе не исследована.

\textit{Ограниченный масштаб LLM-прогона.} Политика на основе языковой модели запускалась только на двенадцати точках гиперкуба; LLM-агентский симулятор работал на тех же двенадцати точках с одним случайным зерном. Использована экономичная модель из соображений бюджета --- это сознательное ограничение и направление развития работы.

\textit{Не полностью воспроизводимые ответы LLM-агентов.} Параметрическая часть эксперимента воспроизводима полностью. Ответы LLM-агентов не кэшируются: повторный прогон LLM-симулятора может дать другой набор индивидуальных решений из-за стохастичности языковой модели. Метаданные условий запуска фиксируются.

\textit{Конфаундированность анализа по оси индекса варианта программы.} Точки гиперкуба с разными уровнями индекса варианта программы не сопоставлены между собой по остальным осям. Агрегаты по этой оси (см. таблицу~\ref{tab:program-effect}) --- диагностические, не причинные.

\textit{Узкая выборка для согласованности по метрикам переполнения.} Корреляция Спирмена между двумя симуляторами по метрикам, непосредственно связанным с переполнением, опирается на две невырожденные точки из двенадцати (см. таблицу~\ref{tab:cross-validation}). Вывод о согласованности по этим метрикам остаётся диагностическим до прогонов на конфигурациях с большей долей risk-positive точек.

\textit{Сценарная, а не прогнозная интерпретация.} Полученные значения трактуются как сравнительные внутри принятой модели; утверждения вида <<при политике P фактическое превышение составит N процентов>> в работе не делаются и из результатов не следуют.


\section{Выводы по главе}
\label{sec:4-9}

Содержательный итог экспериментальной валидации программного комплекса формулируется в семи тезисах.

\begin{enumerate}
\item \textbf{Граничная верификация пройдена полностью.} Все четыре свойства, обязательные для корректной модели рассматриваемого класса, на текущей реализации выполняются; 10 из 10 автоматизированных тестов завершились без отказов (таблица~\ref{tab:ec-verification}).

\item \textbf{Политика с учётом загрузки залов снижает средний риск при сохранении релевантности.} В полном прогоне она даёт меньшее среднее относительное превышение вместимости (0.0336), чем контрольная политика (0.0395) и политика по релевантности (0.0449), при практически неизменной средней релевантности --- различие в пределах 0.002 (таблица~\ref{tab:central-metrics}).

\item \textbf{В попарном сравнении политика по релевантности не выигрывает у политики с учётом загрузки залов ни в одной из 50 точек гиперкуба} ни строго (\texttt{win\_strict = 0}), ни за пределами допуска $\varepsilon = 0.005$ (\texttt{win\_eps = 0}). Обратное направление выигрыша зафиксировано в 26~\% точек строго и в 22~\% точек за пределами допуска (таблица~\ref{tab:pairwise-overload}).

\item \textbf{На risk-positive подмножестве политика с учётом загрузки залов устойчиво лидирует.} В 13 точках с ненулевым риском хотя бы у одной из политик она не уступает максимуму из двух остальных в 13 / 13 точек в пределах допуска и строго снижает риск в 11 / 13 точек (таблицы~\ref{tab:risk-positive} и~\ref{tab:risk-reductions}).

\item \textbf{Trade-off между снижением риска и потерей релевантности редок.} На полном прогоне он встречается в 7.3~\% комбинаций, на подмножестве с участием LLM-политики --- в 6.3~\% (таблица~\ref{tab:tradeoff}).

\item \textbf{Согласованность параметрического и LLM-агентского симуляторов наблюдается по части показателей.} По средней релевантности --- медиана $\rho = 0.80$ на 12 точках и совпадение верхнего лидера на 11 / 12 точках. По дисперсии загрузки залов --- совпадение верхнего лидера на 11 / 12 точках при медиане $\rho = 0.40$ ниже формального порога. По метрикам переполнения параметрическое ранжирование вырождается на 10 из 12 точек, и Spearman $\rho$ оценивается на 2 невырожденных точках; вывод по этим метрикам остаётся диагностическим (таблица~\ref{tab:cross-validation}). Эти результаты частично поддерживают третье защищаемое положение: реализованы два механизма отклика и проведена перекрёстная проверка, но по метрикам переполнения согласованность остаётся диагностической.

\item \textbf{Результаты --- сценарное сравнительное утверждение, не прогноз.} Полученные значения относятся к сравнительному анализу политик и вариантов программы в рамках принятой модели на программе Mobius 2025 Autumn. Они не утверждают воспроизведения фактического поведения аудитории, не выделяют единственной оптимальной политики рекомендаций и не претендуют на универсальность для произвольной программы конференции.
\end{enumerate}

Содержательное значение полученного отчёта для организатора программы конференции состоит в следующем: использование сравнительного отчёта позволяет заранее идентифицировать сценарии параметрической неопределённости, в которых программа уязвима к перегрузке залов, и принять решение о выборе политики рекомендаций или о структурной модификации программы до проведения мероприятия --- на основании сравнительной устойчивости показателей риска внутри принятой модели, а не на основании прогноза фактической посещаемости.

Полное содержание сформированного сравнительного отчёта --- карта ожидаемой загрузки залов, перечень горячих точек программы, попарные сравнения политик, диагностика по уровню индекса варианта программы, анализ чувствительности по параметрическим осям и блок согласованности двух симуляторов --- сохранено в составе результатов и предъявляется как итоговый артефакт системы поддержки принятия решений для организатора программы конференции.
